% Diego Sebastián Aliaga Damián
% 2024200481M
% Tema 96: Fintech y financiamiento participativo (crowdfunding) en el Perú: oferta y condiciones
% Artículo científico (journal), Finanzas I (055D), Unidad II, UNCP, 2026-II
%
% Compilación (desde esta carpeta):  latexmk -pdf articulo.tex
% Antes de compilar:                 python ..\herramientas\articulo\generar_insumos_articulo.py
% Las tablas y figuras se leen de ../salidas (generadas por codigo/04_analisis.py) y las cifras del
% texto de generado/cifras.tex (derivadas de salidas/resultados_04.json): nada se transcribe a mano.

\documentclass[12pt]{article}

% ---------- Página y tipografía (APA 7: márgenes de 2,54 cm, Times 12, doble espacio) ----------
\usepackage[a4paper,margin=2.54cm]{geometry}
\usepackage[T1]{fontenc}
\usepackage[utf8]{inputenc}
\usepackage[spanish,es-noshorthands,es-tabla,es-nodecimaldot]{babel}
\usepackage{newtxtext,newtxmath}
\usepackage{setspace}
\doublespacing
\setlength{\parindent}{1.27cm}
\usepackage{csquotes}
\usepackage[expansion=false]{microtype}  % sin expansión de fuentes: compila igual en cualquier instalación

% ---------- Encabezado: número de página arriba a la derecha ----------
\usepackage{fancyhdr}
\pagestyle{fancy}
\fancyhf{}
\fancyhead[R]{\thepage}
\renewcommand{\headrulewidth}{0pt}

% ---------- Títulos de sección según APA 7 (niveles 1 a 3) ----------
\usepackage{titlesec}
\titleformat{\section}{\centering\bfseries\normalsize}{}{0pt}{}
\titleformat{\subsection}{\bfseries\normalsize}{}{0pt}{}
\titleformat{\subsubsection}{\bfseries\itshape\normalsize}{}{0pt}{}
\titlespacing*{\section}{0pt}{12pt}{0pt}
\titlespacing*{\subsection}{0pt}{12pt}{0pt}
\titlespacing*{\subsubsection}{0pt}{12pt}{0pt}

% ---------- Tablas y figuras con formato APA 7 ----------
\usepackage{graphicx,booktabs,array,float,pdflscape,ragged2e,longtable,tabularx}  % 04 escribe las tablas con tabularx
\usepackage{newunicodechar}  % símbolos que aparecen en las tablas generadas por 04
\newunicodechar{≤}{\ensuremath{\leq}}
\newunicodechar{≥}{\ensuremath{\geq}}
\newunicodechar{Δ}{\ensuremath{\Delta}}
\newunicodechar{≈}{\ensuremath{\approx}}
\usepackage{caption}
\captionsetup{labelfont=bf,textfont=it,labelsep=newline,justification=RaggedRight,singlelinecheck=false,
              skip=4pt,font={normalsize,singlespacing}}
\captionsetup[table]{position=top}
\captionsetup[figure]{position=top}
\graphicspath{{../salidas/figuras/}}
\newcommand{\tablas}{../salidas/tablas}
\usepackage{etoolbox}
\AtBeginEnvironment{table}{\singlespacing\small}  % APA 7 admite interlineado sencillo dentro de tablas y figuras
\AtBeginEnvironment{figure}{\singlespacing}
\newcommand{\notaapa}[1]{\par\vspace{4pt}{\small\RaggedRight\textit{Nota.} #1\par}}

% ---------- Referencias en APA 7 (biblatex-apa + biber) ----------
\usepackage[style=apa,backend=biber]{biblatex}
\addbibresource{referencias.bib}
\usepackage{xurl}
\usepackage[hidelinks]{hyperref}

% ---------- Cifras del texto generadas desde salidas/resultados_04.json ----------
% Generado por herramientas/articulo/generar_insumos_articulo.py; no editar a mano.
\newcommand{\cifra}[1]{\ifcsname cifra@#1\endcsname\csname cifra@#1\endcsname\else\textbf{??#1??}\fi}
\expandafter\def\csname cifra@acfPrevia\endcsname{0.84}
\expandafter\def\csname cifra@bootstrap\endcsname{10\,000}
\expandafter\def\csname cifra@coberturaPct\endcsname{90.3}
\expandafter\def\csname cifra@colocacion2022\endcsname{54.5}
\expandafter\def\csname cifra@colocacion2023\endcsname{77.9}
\expandafter\def\csname cifra@colocacion2024\endcsname{83.4}
\expandafter\def\csname cifra@colocacion2025\endcsname{88.5}
\expandafter\def\csname cifra@colocacion2026\endcsname{94.2}
\expandafter\def\csname cifra@colocacionEmp\endcsname{89.5}
\expandafter\def\csname cifra@colocacionPer\endcsname{51.4}
\expandafter\def\csname cifra@columnasPanel\endcsname{35}
\expandafter\def\csname cifra@costoAcumulada\endcsname{0.76}
\expandafter\def\csname cifra@costoMediana\endcsname{5.27}
\expandafter\def\csname cifra@costoTrimestres\endcsname{17}
\expandafter\def\csname cifra@cotaMaxAnual\endcsname{0.22}
\expandafter\def\csname cifra@cotaPctMax\endcsname{1.7}
\expandafter\def\csname cifra@cotaPctMaxDesde2024\endcsname{0.2}
\expandafter\def\csname cifra@crecMensualPrevio\endcsname{8.9}
\expandafter\def\csname cifra@eUnoN\endcsname{8}
\expandafter\def\csname cifra@eUnoPMin\endcsname{0.116}
\expandafter\def\csname cifra@eUnoRhoMax\endcsname{0.01}
\expandafter\def\csname cifra@eUnoRhoMin\endcsname{-0.22}
\expandafter\def\csname cifra@eeffEscala\endcsname{8}
\expandafter\def\csname cifra@eeffEscalaCorregida\endcsname{3}
\expandafter\def\csname cifra@eeffEscalaExcluida\endcsname{5}
\expandafter\def\csname cifra@fechaCorte\endcsname{septiembre de 2026}
\expandafter\def\csname cifra@fechaInicio\endcsname{enero de 2021}
\expandafter\def\csname cifra@filasPanel\endcsname{414}
\expandafter\def\csname cifra@hDosFisherOR\endcsname{2.84}
\expandafter\def\csname cifra@hDosFisherP\endcsname{0.127}
\expandafter\def\csname cifra@hDosP\endcsname{< 0.001}
\expandafter\def\csname cifra@hDosZ\endcsname{14.88}
\expandafter\def\csname cifra@hUnoICinf\endcsname{5.63}
\expandafter\def\csname cifra@hUnoICsup\endcsname{11.27}
\expandafter\def\csname cifra@hUnoMedianaAntes\endcsname{3.32}
\expandafter\def\csname cifra@hUnoMedianaDespues\endcsname{11.37}
\expandafter\def\csname cifra@hUnoNAntes\endcsname{42}
\expandafter\def\csname cifra@hUnoNDespues\endcsname{10}
\expandafter\def\csname cifra@hUnoP\endcsname{< 0.001}
\expandafter\def\csname cifra@hUnoPCota\endcsname{< 0.001}
\expandafter\def\csname cifra@hUnoPSinLider\endcsname{0.124}
\expandafter\def\csname cifra@hUnoU\endcsname{387}
\expandafter\def\csname cifra@hhi2022\endcsname{8\,882}
\expandafter\def\csname cifra@hhi2023\endcsname{9\,770}
\expandafter\def\csname cifra@hhi2024\endcsname{9\,999}
\expandafter\def\csname cifra@hhi2025\endcsname{9\,977}
\expandafter\def\csname cifra@hhi2026\endcsname{9\,321}
\expandafter\def\csname cifra@inicioFPF\endcsname{junio de 2022}
\expandafter\def\csname cifra@interanualAntes\endcsname{169}
\expandafter\def\csname cifra@interanualDespues\endcsname{27}
\expandafter\def\csname cifra@interanualP\endcsname{1.000}
\expandafter\def\csname cifra@lider\endcsname{Inversiones.io}
\expandafter\def\csname cifra@mesesActividadFPF\endcsname{52}
\expandafter\def\csname cifra@mesesDespues\endcsname{10}
\expandafter\def\csname cifra@mesesObtenidos\endcsname{374}
\expandafter\def\csname cifra@mesesSobreTendencia\endcsname{0}
\expandafter\def\csname cifra@monto2022\endcsname{6.46}
\expandafter\def\csname cifra@monto2023\endcsname{16.24}
\expandafter\def\csname cifra@monto2024\endcsname{54.28}
\expandafter\def\csname cifra@monto2025\endcsname{120.28}
\expandafter\def\csname cifra@monto2026\endcsname{104.71}
\expandafter\def\csname cifra@montoMedio2022\endcsname{58.8}
\expandafter\def\csname cifra@montoMedio2023\endcsname{115.2}
\expandafter\def\csname cifra@montoMedio2024\endcsname{215.4}
\expandafter\def\csname cifra@montoMedio2025\endcsname{278.4}
\expandafter\def\csname cifra@montoMedio2026\endcsname{320.2}
\expandafter\def\csname cifra@montoTotal\endcsname{301.97}
\expandafter\def\csname cifra@noObtenidos\endcsname{40}
\expandafter\def\csname cifra@noObtenidosVentana\endcsname{28}
\expandafter\def\csname cifra@participacionLider\endcsname{98.5}
\expandafter\def\csname cifra@pctEmpresarialPrincipal\endcsname{99.5}
\expandafter\def\csname cifra@pctPersonalPrincipal\endcsname{95.3}
\expandafter\def\csname cifra@perLiderFin\endcsname{9}
\expandafter\def\csname cifra@perLiderPub\endcsname{12}
\expandafter\def\csname cifra@placeboN\endcsname{5}
\expandafter\def\csname cifra@placeboPMax\endcsname{< 0.001}
\expandafter\def\csname cifra@plataformasAutorizadas\endcsname{6}
\expandafter\def\csname cifra@plataformasConFinanciamiento\endcsname{3}
\expandafter\def\csname cifra@plataformasVigentes\endcsname{3}
\expandafter\def\csname cifra@principalPersonal\endcsname{Afluenta}
\expandafter\def\csname cifra@proyFin2022\endcsname{110}
\expandafter\def\csname cifra@proyFin2023\endcsname{141}
\expandafter\def\csname cifra@proyFin2024\endcsname{252}
\expandafter\def\csname cifra@proyFin2025\endcsname{432}
\expandafter\def\csname cifra@proyFin2026\endcsname{327}
\expandafter\def\csname cifra@proyFinTotal\endcsname{1\,262}
\expandafter\def\csname cifra@rDosTendencia\endcsname{0.81}
\expandafter\def\csname cifra@saltoTEA\endcsname{19.2}
\expandafter\def\csname cifra@semilla\endcsname{0481}
\expandafter\def\csname cifra@solicitudesDiagnostico\endcsname{26}
\expandafter\def\csname cifra@solicitudesHTTP\endcsname{1\,702}
\expandafter\def\csname cifra@solicitudesHTTPok\endcsname{1\,702}
\expandafter\def\csname cifra@varMonto2023\endcsname{151.3}
\expandafter\def\csname cifra@varMonto2024\endcsname{234.3}
\expandafter\def\csname cifra@varMonto2025\endcsname{121.6}
\expandafter\def\csname cifra@varMonto2026\endcsname{14.7}


% ---------- Marca de borrador: debe desaparecer antes de entregar ----------
\usepackage{xcolor}
\newcommand{\pendiente}[1]{\textcolor{red}{[PENDIENTE: #1]}}

\begin{document}

% Portada de trabajo estudiantil (APA 7) con título en español e inglés (consigna 2.2, ítems 1 y 2).
% APA 7 numera también la portada (página 1, arriba a la derecha).
\thispagestyle{fancy}
\begin{center}
\vspace*{3cm}
{\bfseries Fintech y financiamiento participativo (crowdfunding) en el Perú: oferta y condiciones, 2022-2026\par}
\vspace{12pt}
{\bfseries\itshape Fintech and Crowdfunding in Peru: Supply and Financing Terms, 2022-2026\par}
\vspace{2cm}
Diego Sebastián Aliaga Damián\par
Código de matrícula 2024200481M\par
\vspace{12pt}
Facultad de Economía, Universidad Nacional del Centro del Perú\par
Finanzas I (055D), Unidad II, tema N.\textsuperscript{o}~96\par
Dr. Ciro Iván Machacuay Meza\par
6 de diciembre de 2026\par
\end{center}
\clearpage

% Resumen y palabras clave; Abstract y keywords (consigna 2.2, ítems 3 y 4). Se redactan al final.
\section{Resumen}
\noindent\pendiente{redactar al final, con las cifras ya cerradas (150 a 250 palabras)}

\noindent\textit{Palabras clave:} financiamiento participativo financiero, crowdfunding, fintech, crédito a pymes, Perú

\clearpage
\section{Abstract}
\noindent\pendiente{traducción del resumen}

\noindent\textit{Keywords:} lending-based crowdfunding, fintech, small business credit, Peru

\clearpage

% Introducción (consigna 2.2, ítem 5)
\section{Introducción}
\pendiente{sección 6 del orden de redacción}

% Antecedentes (consigna 2.2, ítem 6): estado del arte 2023-2026 y marco normativo
\section{Antecedentes}
\pendiente{sección 3 del orden de redacción}

% Materiales y métodos (consigna 2.2, ítem 7), con la subsección de fuentes de datos y procedimiento de extracción.
% BORRADOR ASISTIDO POR IA (2026-11-11, corregido tras la auditoría metodológica). El estudiante debe reescribirlo
% con sus palabras antes de entregar (regla 2.1: más del 70 % de texto de IA sin humanizar = nota cero).
% Las cifras usan \cifra{...} y no se tocan.
\section{Materiales y métodos}
\pendiente{borrador asistido por IA: reescribir con tus palabras y borrar esta marca}

El estudio tiene un enfoque cuantitativo, no experimental y de alcance descriptivo-comparativo. Describe la oferta de financiamiento participativo financiero (FPF) en el Perú con datos oficiales y la contrasta con las condiciones del crédito tradicional en segmentos comparables. No busca estimar efectos causales: las pruebas estadísticas describen diferencias y asociaciones, sin atribuirlas a una causa.

\subsection{Diseño, población y periodo}

La unidad de análisis es la plataforma-mes. La población está formada por las \cifra{plataformasAutorizadas} sociedades administradoras de plataformas de FPF que registran información ante la Superintendencia del Mercado de Valores (SMV) durante el periodo, \cifra{plataformasVigentes} de ellas vigentes al cierre. Como se trabaja con todas, no hay muestreo: es un censo de la oferta regulada, con los faltantes que se describen más adelante. La actividad está normada por el Título IV del Decreto de Urgencia N.\textsuperscript{o}~013-2020 \parencite{du0132020} y por el reglamento aprobado mediante la Resolución de Superintendente N.\textsuperscript{o}~045-2021-SMV/02 \parencite{smv0452021}, modificado en noviembre de 2025 \parencite{smv0192025}.

El periodo de consulta va de \cifra{fechaInicio} a \cifra{fechaCorte}, es decir, 69 periodos mensuales. Ambos límites se declararon como constantes en el código para que la consulta quede congelada. El primer mes con proyectos financiados registrados en la SMV es \cifra{inicioFPF}, de modo que el FPF tiene \cifra{mesesActividadFPF} meses con actividad. Los meses previos no son relleno: el portal respondió que no había proyectos. Se conservan en la base para mantener el panel balanceado y para que las series del crédito tradicional cubran todo el periodo.

\subsection{Fuentes de datos y procedimiento de extracción}

La base se construyó con dos vías automatizadas, como exige la consigna de la Unidad II: el consumo de interfaces de programación (API) y el rastreo programático (\textit{scraping}) de un portal institucional. Todo el procedimiento está en los scripts del repositorio (Anexo), escritos en Python.

\subsubsection{Vía 1: interfaces de programación}

La primera fuente es BCRPData, la API pública del Banco Central de Reserva del Perú \parencite{bcrpdata2026}. De ella se descargaron diez series mensuales en moneda nacional: la tasa activa promedio de las empresas bancarias (TAMN); las tasas efectivas anuales de los préstamos a pequeña y microempresa de bancos, cajas municipales y cajas rurales, en los tramos de hasta 360 días y de más de 360 días; las de los préstamos de consumo de los bancos, en los mismos dos tramos, y el índice de precios al consumidor de Lima Metropolitana (base diciembre de 2021 = 100). El script no confía solo en el código de cada serie: compara el nombre oficial que devuelve la API con palabras clave esperadas y rechaza la respuesta si alguna no coincide.

La segunda fuente es la API OData de información financiera de la SMV \parencite{smvodata2026}, que publica las principales cuentas de los estados financieros de las empresas supervisadas. Se consultaron los 23 trimestres del periodo y se conservaron todas las filas recibidas; las sociedades administradoras de plataformas de FPF se identifican después, en la etapa de limpieza. Antes de guardar cada respuesta se comprueba que esté completa y que tenga las columnas necesarias, y una respuesta vacía nunca reemplaza un archivo que ya tenía datos.

\subsubsection{Vía 2: rastreo del portal de la SMV}

La información sobre la actividad de las plataformas proviene de la página «Información de Proyectos» del Sistema de Información del Mercado de Valores \parencite{smvsimv2026}. Esta página no ofrece archivos descargables: muestra una tabla después de que el usuario elige la plataforma, el año y el mes en un formulario. El script reproduce ese procedimiento. Primero solicita la página y recupera los campos ocultos que el formulario exige devolver; luego envía la consulta para cada plataforma y cada mes del periodo, y lee la tabla de resultados, que informa por modalidad y tipo de receptor (empresarial o personal) los proyectos publicados, los proyectos financiados y el monto total financiado.

Una respuesta solo se guarda si contiene la tabla de resultados, aunque esta diga que no hay proyectos, y si el formulario devuelto confirma la plataforma, el año y el mes consultados; de lo contrario se reintenta con un formulario nuevo. Entre solicitudes se respeta una pausa de tres segundos y todas se identifican con un agente de usuario del proyecto. El sitio de la SMV no publica un archivo \texttt{robots.txt} propio, por lo que se aplicó una política conservadora y se guardó la evidencia de esa consulta.

Las consultas sobre las plataformas que ya no operan presentaron una falla intermitente del portal: el formulario confirmaba la consulta, pero no mostraba la tabla, aunque en el caso que se revisó manualmente el dato sí aparecía en un navegador. Tras reintentar con formularios nuevos, se obtuvieron \cifra{mesesObtenidos} de los \cifra{filasPanel} plataforma-mes posibles (\cifra{coberturaPct}~\%); los \cifra{noObtenidos} restantes corresponden a plataformas retiradas y quedan marcados como no obtenidos. La incidencia se documentó y se informó al docente, que estuvo de acuerdo con tratarla sin sustituir la fuente.

\subsubsection{Integridad y verificación}

Las respuestas originales se guardaron sin editar, y el hash SHA-256 de cada una quedó registrado en el log de ejecución junto con la fecha, la hora y el código HTTP. Las corridas de extracción del 7 y el 8 de noviembre de 2026 hicieron \cifra{solicitudesHTTP} solicitudes, y las \cifra{solicitudesHTTPok} respondieron con el código 200; el log registra además \cifra{solicitudesDiagnostico} solicitudes de pruebas de diagnóstico del portal. Como control externo, los totales anuales construidos con el portal se compararon con las cifras de la SMV que recogió la prensa \parencite{ramirez2026}: para 2025, la base registra \cifra{proyFin2025} proyectos financiados por S/~\cifra{monto2025} millones, frente a 432 proyectos y S/~120.3 millones. Una auditoría automatizada, incluida en el repositorio, vuelve a generar la base y las salidas desde los archivos originales y coteja una muestra aleatoria de filas con las fuentes en vivo; en su ejecución del 11 de noviembre de 2026, todos los valores cotejados con BCRPData y con el portal coincidieron.

\subsection{Limpieza e integración}

Un tercer script lee únicamente los archivos crudos. Convierte los conteos y montos a números, reorganiza la tabla del portal en columnas por tipo de receptor y usa la fila de totales que publica la SMV solo como control: en todas las consultas con registros, el total coincide con la suma de sus componentes. Los estados financieros, que son trimestrales, se asignan al mes de cierre de cada trimestre sin interpolar, y las tasas del BCRP se asignan a su mes. El portal y la API de la SMV se unen por el RUC de la sociedad administradora y el periodo; el cruce es exacto porque ambas fuentes, de la misma institución, usan la misma razón social. Las series del BCRP, que describen a todo el sistema y no a cada plataforma, se unen solo por el periodo. El resultado es un panel balanceado de \cifra{filasPanel} filas y \cifra{columnasPanel} columnas.

No se imputaron datos faltantes. Una variable de control distingue tres casos: los meses con registros, los meses en que el portal respondió que no había proyectos, que el análisis cuenta como cero, y los meses no obtenidos, que se dejan vacíos. En consecuencia, los agregados del FPF son una cota inferior. Para medir cuánto podría cambiar el resultado, se asignó a cada mes no obtenido el mayor monto mensual que su plataforma registró en todo el periodo. Con esa cota, de los \cifra{noObtenidos} meses no obtenidos, \cifra{noObtenidosVentana} caen dentro del periodo con actividad, y el monto anual podría aumentar como máximo en S/~\cifra{cotaMaxAnual} millones (\cifra{cotaPctMax}~\% del monto del año más afectado y no más de \cifra{cotaPctMaxDesde2024}~\% desde 2024).

Los valores atípicos se identificaron con la regla de Tukey (tres veces el rango intercuartílico) y se conservaron, por tratarse de datos oficiales. La revisión detectó además \cifra{eeffEscala} registros de estados financieros con un activo superior a S/~50 millones expresado en miles, una escala inverosímil para estas empresas: todo indica que se reportaron en soles y no en miles de soles. En la base se mantiene el valor publicado, y en el análisis se distinguen dos casos. En \cifra{eeffEscalaCorregida} registros, la misma empresa tiene trimestres con la escala correcta y la división entre mil vuelve continuas sus series de activos e ingresos; esos registros se dividen entre mil. En los otros \cifra{eeffEscalaExcluida}, todos de una empresa sin trimestres de referencia, no hay forma interna de comprobar la escala: sus saldos de activo, pasivo y patrimonio no se usan, y sus ingresos y utilidad, que son cero, se conservan.

\subsection{Variables e indicadores}

Las variables provienen de las tres fuentes y se definen en el diccionario del Anexo. A partir de ellas se calcularon los indicadores del análisis:

\begin{itemize}
\item \textit{Concentración de la oferta}, con el índice de Herfindahl-Hirschman sobre el monto financiado, $\mathit{IHH}_t=\sum_i (100\,s_{it})^2$, donde $s_{it}$ es la participación de la plataforma $i$ en el año $t$.
\item \textit{Tasa de colocación}, como cociente entre proyectos financiados y publicados. Se calcula sobre sumas anuales, de doce meses o de subperiodos, porque un proyecto puede financiarse en un mes posterior al de su publicación y el cociente mensual puede superar la unidad.
\item \textit{Monto medio por proyecto financiado}, por periodo y tipo de receptor.
\item \textit{Ingresos de la plataforma sobre el monto financiado}, como aproximación al costo de intermediación: ingresos del trimestre multiplicados por mil y divididos entre el monto financiado en ese trimestre. Solo se calcula en trimestres con los tres meses del portal obtenidos y, en el cálculo anual, con al menos S/~1 millón financiado. No es un precio para el receptor: mide cuánto retiene la plataforma por cada sol que intermedia.
\item \textit{Rentabilidad sobre activos} (ROA) de las sociedades administradoras, indicador que la literatura sobre eficiencia de las plataformas de crédito considera central \parencite{pampurini2024,onorato2024}. Se calcula como $\mathit{ROA}=100\cdot\frac{4}{n}\sum_q U_q \big/ \bar{A}$, donde $U_q$ es la utilidad neta de cada uno de los $n$ trimestres disponibles del año y $\bar{A}$ es el activo medio; el factor $4/n$ anualiza los años incompletos.
\end{itemize}

Los dos últimos indicadores dependen de un supuesto: que la SMV publica los ingresos y la utilidad de cada trimestre, no los acumulados del año. El diccionario de la fuente no lo precisa, pero en las dos plataformas con mayores ingresos los ingresos bajan de un trimestre a otro dentro de un mismo año, algo imposible si fueran acumulados. La lectura acumulada se muestra solo como cota en el análisis de sensibilidad.

\subsection{Estrategia de comparación con el crédito tradicional}

La SMV no publica la tasa ni el plazo de cada operación de FPF, de modo que el precio del financiamiento participativo no se observa. Por eso, la comparación con el crédito tradicional se limita a las dimensiones observables y usa, para cada condición de la consigna, una variable sustituta aceptada por el docente (Tabla~\ref{tab:comparabilidad}). La tasa del crédito tradicional sirve solo como referencia del sistema y no como término de comparación de precios, porque no hay una tasa del FPF con la cual contrastarla. El sector se aproxima con el tipo de receptor: los proyectos empresariales se ubican frente a las tasas de préstamos a pequeña y microempresa, y los personales frente a las de consumo. La TAMN se reporta como referencia general del crédito bancario.

Para el plazo se usa el tramo de hasta 360 días, por dos razones. La evidencia reciente muestra que los prestamistas tecnológicos se especializan en créditos sin garantía para proyectos de corto plazo \parencite{serfes2025}, y la prensa informó plazos promedio de 217 días en 2024 y de 201 días en 2025 para los préstamos de financiamiento participativo \parencite{ramirez2026}. El tramo de más de 360 días se usa como sensibilidad. La serie bancaria de pequeña y microempresa de hasta 360 días sube \cifra{saltoTEA} puntos porcentuales entre febrero y marzo de 2025, sin un movimiento parecido en las cajas ni en el crédito de consumo; como no se encontró una nota metodológica que lo explique, el año 2025 se reporta dividido en dos subperiodos, antes y desde ese mes.

\begin{table}[H]
\caption{Comparabilidad entre el financiamiento participativo financiero y el crédito tradicional}
\label{tab:comparabilidad}
\input{\tablas/t6_comparabilidad.tex}
\notaapa{Elaboración propia a partir del diseño del estudio. La tasa y el plazo del financiamiento participativo no se publican por operación; la tasa efectiva anual del crédito tradicional se usa solo como referencia del sistema.}
\end{table}

\subsection{Métodos estadísticos}

El análisis combina estadística descriptiva con tres contrastes. Para la hipótesis H1, según la cual el monto mensual financiado es mayor después de la modificación del reglamento, se comparan los \cifra{hUnoNAntes} meses previos y los \cifra{hUnoNDespues} meses posteriores, con diciembre de 2025 como primer mes completo de vigencia de la modificación, mediante la prueba U de Mann-Whitney unilateral. La diferencia de medianas se acompaña de un intervalo de confianza al 95~\% obtenido por \textit{bootstrap} de percentiles con \cifra{bootstrap} réplicas. Esta comparación es de niveles y tiene dos limitaciones: supone meses independientes, aunque el monto mensual está muy autocorrelacionado, y no separa la modificación del crecimiento que el mercado ya traía. Por eso se agregan tres controles: la misma prueba con \cifra{placeboN} cortes ficticios ubicados antes de la modificación, que deberían salir no significativos si la prueba detectara solo el cambio normativo; la comparación del crecimiento interanual antes y después del corte, y la posición de los meses posteriores respecto de la tendencia log-lineal estimada con los meses previos.

La hipótesis H2 compara la tasa de colocación de los proyectos empresariales con la de los personales mediante una prueba \textit{z} de diferencia de proporciones, con intervalos de Wilson al 95~\%. La prueba supone proyectos independientes, aunque los proyectos de una misma plataforma y un mismo mes pueden no serlo, por lo que los intervalos deben leerse como aproximados. Además, cada tipo de receptor se concentra en una plataforma distinta, de modo que la comparación agregada también refleja diferencias entre modelos de negocio. Para aislar el tipo de receptor, se repite el contraste dentro de la plataforma líder con la prueba exacta de Fisher, adecuada cuando uno de los grupos es pequeño.

La relación exploratoria E1 examina si los cambios mensuales del monto financiado se asocian con los cambios de las tasas tradicionales. Para no confundir tendencias comunes con asociación, se usan primeras diferencias y la correlación de Spearman, con intervalos obtenidos por \textit{bootstrap} de bloques móviles de tres meses. Se evalúan \cifra{eUnoN} comparadores: pequeña y microempresa de bancos, cajas municipales y cajas rurales, y consumo bancario, cada uno en sus dos tramos de plazo. Con algo más de cincuenta cambios mensuales, la prueba tiene poca potencia, por lo que un resultado no significativo no equivale a ausencia de asociación.

La robustez de H1 se evaluó, además, con el corte en noviembre de 2025, en soles constantes de diciembre de 2021, con la cota superior de los meses no obtenidos, solo con la plataforma líder y sin ella. Todos los procedimientos aleatorios usan la semilla \cifra{semilla}, formada por los cuatro últimos dígitos de la matrícula del autor. Las tablas y figuras de la sección siguiente se generan con el script de análisis y se insertan en el artículo sin transcripción manual.

\subsection{Consideraciones éticas}

Solo se usó información pública y agregada; la base no contiene datos personales protegidos por la Ley N.\textsuperscript{o}~29733. Las fuentes que aplican protecciones contra el acceso automatizado, como el portal de estadísticas de la Superintendencia de Banca, Seguros y AFP, se descartaron en lugar de eludirlas. Las incidencias de la extracción, incluidas las del propio autor durante los sondeos iniciales, están documentadas en el repositorio, y el uso de herramientas de inteligencia artificial se declara al final del artículo.

% Resultados (consigna 2.2, ítem 8): tablas y figuras importadas de ../salidas
\section{Resultados}
\pendiente{sección 2 del orden de redacción}

% Discusión (consigna 2.2, ítem 9)
\section{Discusión}
\pendiente{sección 4 del orden de redacción}

% Conclusiones (consigna 2.2, ítem 10). BORRADOR ASISTIDO POR IA (2026-11-12): reescribir con tus palabras.
\section{Conclusiones}

La oferta de financiamiento participativo financiero en el Perú creció con fuerza entre 2022 y 2026, de S/~\cifra{monto2022} a S/~\cifra{monto2025} millones anuales, pero está muy concentrada: solo \cifra{plataformasConFinanciamiento} de \cifra{plataformasAutorizadas} administradoras financiaron proyectos y una reúne el \cifra{participacionLider}~\% del monto. Es, sobre todo, un mercado de préstamos a empresas con montos crecientes, en el que solo la plataforma líder es rentable y retiene alrededor de \cifra{costoMediana}~\% de lo que financia. Las diferencias de colocación entre receptores reflejan el modelo de negocio de cada plataforma, y el mayor monto tras la modificación de 2025 es coherente con el crecimiento previo, por lo que no puede atribuirse a la norma.

Frente a un crédito tradicional de tasas altas, el FPF no puede compararse en precio: mientras no se publiquen la tasa y el plazo de sus operaciones, no podrá saberse si ofrece mejores condiciones que la banca y las cajas.

% Declaración de uso de inteligencia artificial (consigna 2.1, «Sobre el uso de IA»; rúbrica C6)
\section{Declaración de uso de inteligencia artificial}
\pendiente{el estudiante la redacta con sus palabras: qué partes se apoyaron en IA, cómo verificó los resultados y qué decidió por su cuenta}


\clearpage
\section{Referencias}
\begingroup
\setlength{\parindent}{0pt}
\printbibliography[heading=none]
\endgroup

\clearpage
% Anexo de datos y código (consigna 2.2, ítem 12): repositorio, endpoints, diccionario, fecha de corte y hash.
% Todas las tablas de este anexo se generan desde los archivos del proyecto (articulo/generado/).
\section{Anexo. Datos y código}

\subsection{Repositorio, fechas y verificación}
\begin{singlespace}
% Generado por herramientas/articulo/generar_insumos_articulo.py; no editar a mano.
\begin{longtable}{@{}>{\RaggedRight\arraybackslash}p{3.8cm}>{\RaggedRight\arraybackslash}p{11.3cm}@{}}
\toprule
Elemento & Detalle \\
\midrule
\endhead
Repositorio & \url{https://github.com/Diego-Aliaga-Damian/finanzas1-u2-crowdfunding-peru} \\
Commit de referencia & \texttt{0a69dc1} \\
Periodo consultado & 2021-01 a 2026-09 (constantes FECHA\_INICIO y FECHA\_CORTE) \\
Fechas de extracción & 2026-11-07, 2026-11-08 (según log\_ejecucion.txt) \\
Archivo procesado & \texttt{datos\_procesados/datos\_procesados\_2024200481M.csv}: 414 filas, 35 columnas \\
SHA-256 del archivo procesado & \texttt{cd48b21f22844cff650f5257a324b632}\newline\texttt{d5afac7a009dbaceebe7456e9a7e1b8c} \\
Orden de ejecución & \texttt{01\_extraccion\_api.py} $\rightarrow$ \texttt{02\_scraping\_web.py} $\rightarrow$ \texttt{03\_limpieza\_datos.py} $\rightarrow$ \texttt{04\_analisis.py} \\
Verificación & \texttt{python -m pytest} y \texttt{python codigo/auditoria.py} (con \texttt{-{}-en-vivo} para cotejar con las fuentes) \\
Entorno & Python 3.14.7; requests 2.34.2, beautifulsoup4 4.15.0, pandas 3.1.1, numpy 2.5.4, scipy 1.18.1, matplotlib 3.11.2, pandera 0.35.0, pytest 9.2.0 \\
\bottomrule
\end{longtable}

\end{singlespace}

\subsection{Fuentes, endpoints y parámetros de extracción}
\begin{singlespace}
% Generado por herramientas/articulo/generar_insumos_articulo.py; no editar a mano.
\begin{longtable}{@{}>{\RaggedRight\arraybackslash}p{1.6cm}>{\RaggedRight\arraybackslash}p{2.8cm}>{\RaggedRight\arraybackslash}p{5.8cm}>{\RaggedRight\arraybackslash}p{3.9cm}@{}}
\toprule
Vía & Fuente & Endpoint o URL & Parámetros \\
\midrule
\endhead
API & BCRP, BCRPData & \url{https://estadisticas.bcrp.gob.pe/estadisticas/series/api/{codigos}/json/{inicio}/{fin}} & 10 series mensuales; 2021-1 a 2026-9 \\
API & SMV, OData PrincipalCuentaSMV & \url{https://mvnet.smv.gob.pe/SMV.OData.Api/api/InfFinanciera/PrincipalCuentaSMV} & sEjercicio=2021-2026; sPeriodo=1-4; sTipoInf=I (23 trimestres) \\
Scraping & SMV, SIMV «Información de Proyectos» & \url{https://www.smv.gob.pe/SIMV/Frm_ProyectosFPF?data=785FE29A0201FCD759A08F94A0F74DBB4FD1B4D805} & POST ASP.NET por plataforma, año y mes; pausa de 3 s; 6 plataformas × 69 meses \\
\bottomrule
\end{longtable}

\end{singlespace}

\subsection{Series del BCRP}
\begin{singlespace}
% Generado por herramientas/articulo/generar_insumos_articulo.py; no editar a mano.
{\small
\begin{longtable}{@{}>{\RaggedRight\arraybackslash}p{2.1cm}>{\RaggedRight\arraybackslash}p{4.8cm}>{\RaggedRight\arraybackslash}p{8.0cm}@{}}
\toprule
Código & Variable & Nombre oficial en BCRPData \\
\midrule
\endhead
PN07807NM & \texttt{tamn} & Tasas de interés activas y pasivas promedio de las empresas bancarias en MN (términos efectivos anuales) - Activas - TAMN \\
PN07843NM & \texttt{tea\_bcos\_pyme\_360} & Tasas de interés activas promedio de las empresas bancarias por modalidad (términos efectivos anuales) - MN - Pequeña y Microempresa - Préstamos - Hasta 360 días \\
PN07844NM & \texttt{tea\_bcos\_pyme\_mas360} & Tasas de interés activas promedio de las empresas bancarias por modalidad (términos efectivos anuales) - MN - Pequeña y Microempresa - Préstamos - Más 360 días \\
PN07846NM & \texttt{tea\_bcos\_consumo\_360} & Tasas de interés activas promedio de las empresas bancarias por modalidad (términos efectivos anuales) - MN - Consumo - Préstamos - Hasta 360 días \\
PN07847NM & \texttt{tea\_bcos\_consumo\_mas360} & Tasas de interés activas promedio de las empresas bancarias por modalidad (términos efectivos anuales) - MN - Consumo - Préstamos - Más 360 días \\
PN07861NM & \texttt{tea\_cmac\_pyme\_360} & Tasas de interés activas promedio de las cajas municipales de ahorro y crédito por modalidad (\% términos efectivos anuales) - MN - Pequeña y Microempresa - Préstamos - Hasta 360 días \\
PN07862NM & \texttt{tea\_cmac\_pyme\_mas360} & Tasas de interés activas promedio de las cajas municipales de ahorro y crédito por modalidad (\% términos efectivos anuales) - MN - Pequeña y Microempresa - Préstamos - Más 360 días \\
PN07879NM & \texttt{tea\_crac\_pyme\_360} & Tasas de interés activas promedio de las cajas rurales de ahorro y crédito por modalidad (\% términos efectivos anuales) - MN - Pequeña y Microempresa - Préstamos - Hasta 360 días \\
PN07880NM & \texttt{tea\_crac\_pyme\_mas360} & Tasas de interés activas promedio de las cajas rurales de ahorro y crédito por modalidad (\% términos efectivos anuales) - MN - Pequeña y Microempresa - Préstamos - Más 360 días \\
PN38705PM & \texttt{ipc} & Índice de precios Lima Metropolitana (índice Dic.2021 = 100) - Índice de Precios al Consumidor (IPC) \\
\bottomrule
\end{longtable}
}

\end{singlespace}

\subsection{Diccionario de variables}
\begin{singlespace}
% Generado por herramientas/articulo/generar_insumos_articulo.py; no editar a mano.
{\footnotesize\setlength{\tabcolsep}{4pt}
\begin{longtable}{@{}>{\RaggedRight\arraybackslash}p{3.8cm}>{\RaggedRight\arraybackslash}p{5.0cm}>{\RaggedRight\arraybackslash}p{2.2cm}>{\RaggedRight\arraybackslash}p{1.7cm}>{\RaggedRight\arraybackslash}p{1.9cm}@{}}
\toprule
Variable & Definición & Unidad & Frecuencia & Fuente \\
\midrule
\endhead
\texttt{periodo} & Mes de referencia (AAAA-MM). Llave común & — & Mensual & Derivada \\
\texttt{anio} & Año del periodo & — & Mensual & Derivada \\
\texttt{mes} & Mes del periodo (1-12) & — & Mensual & Derivada \\
\texttt{ruc} & RUC de la sociedad administradora. Llave común & — & — & SMV-OData \\
\texttt{empresa\_\allowbreak{}id\_\allowbreak{}simv} & Código de la empresa en el formulario del SIMV & — & — & SIMV \\
\texttt{plataforma} & Nombre corto de la plataforma (etiqueta asignada por el autor) & — & — & Derivada \\
\texttt{razon\_\allowbreak{}social} & Razón social de la sociedad administradora & — & — & SIMV \\
\texttt{lista\_\allowbreak{}simv} & Lista del SIMV en que figura: vigentes o retiradas & — & — & SIMV \\
\texttt{estado\_\allowbreak{}simv} & Resultado de la consulta: con\_registros, sin\_registros (el portal no reporta datos) o no\_obtenido (sin respuesta válida; ver incidencias\_fuente.md) & — & Mensual & SIMV \\
\texttt{proy\_\allowbreak{}publicados\_\allowbreak{}emp} & Proyectos de préstamo empresariales publicados en la plataforma en el mes & proyectos & Mensual & SIMV \\
\texttt{proy\_\allowbreak{}financiados\_\allowbreak{}emp} & Proyectos de préstamo empresariales publicados que fueron financiados & proyectos & Mensual & SIMV \\
\texttt{monto\_\allowbreak{}financiado\_\allowbreak{}emp} & Monto total financiado de proyectos de préstamo empresariales & soles corrientes (S/) & Mensual & SIMV \\
\texttt{proy\_\allowbreak{}publicados\_\allowbreak{}per} & Proyectos de préstamo personales publicados en la plataforma en el mes & proyectos & Mensual & SIMV \\
\texttt{proy\_\allowbreak{}financiados\_\allowbreak{}per} & Proyectos de préstamo personales publicados que fueron financiados & proyectos & Mensual & SIMV \\
\texttt{monto\_\allowbreak{}financiado\_\allowbreak{}per} & Monto total financiado de proyectos de préstamo personales & soles corrientes (S/) & Mensual & SIMV \\
\texttt{proy\_\allowbreak{}publicados\_\allowbreak{}valores} & Proyectos publicados en la modalidad valores (todos los receptores) & proyectos & Mensual & SIMV \\
\texttt{proy\_\allowbreak{}financiados\_\allowbreak{}valores} & Proyectos financiados en la modalidad valores & proyectos & Mensual & SIMV \\
\texttt{monto\_\allowbreak{}financiado\_\allowbreak{}valores} & Monto financiado en la modalidad valores & soles corrientes (S/) & Mensual & SIMV \\
\texttt{trimestre\_\allowbreak{}eeff} & Trimestre del estado financiero (AAAA-Tn), asignado a su mes de cierre & — & Trimestral & SMV-OData \\
\texttt{activo\_\allowbreak{}total\_\allowbreak{}miles} & Activo total de la sociedad administradora & miles de S/ & Trimestral & SMV-OData \\
\texttt{pasivo\_\allowbreak{}total\_\allowbreak{}miles} & Pasivo total & miles de S/ & Trimestral & SMV-OData \\
\texttt{patrimonio\_\allowbreak{}total\_\allowbreak{}miles} & Patrimonio total & miles de S/ & Trimestral & SMV-OData \\
\texttt{ingresos\_\allowbreak{}totales\_\allowbreak{}miles} & Ingresos totales del periodo informado & miles de S/ & Trimestral & SMV-OData \\
\texttt{utilidad\_\allowbreak{}neta\_\allowbreak{}miles} & Utilidad neta del periodo informado & miles de S/ & Trimestral & SMV-OData \\
\texttt{fecha\_\allowbreak{}registro\_\allowbreak{}eeff} & Fecha en que la empresa registró el estado financiero en la SMV & fecha (AAAA-MM-DD) & Trimestral & SMV-OData \\
\texttt{tamn} & Tasa activa promedio en moneda nacional (TAMN), empresas bancarias & \% (tasa efectiva anual) & Mensual & PN07807NM \\
\texttt{tea\_\allowbreak{}bcos\_\allowbreak{}pyme\_\allowbreak{}360} & TEA de préstamos a pequeña y microempresa hasta 360 días, empresas bancarias, MN & \% (tasa efectiva anual) & Mensual & PN07843NM \\
\texttt{tea\_\allowbreak{}bcos\_\allowbreak{}pyme\_\allowbreak{}mas360} & TEA de préstamos a pequeña y microempresa a más de 360 días, empresas bancarias, MN & \% (tasa efectiva anual) & Mensual & PN07844NM \\
\texttt{tea\_\allowbreak{}bcos\_\allowbreak{}consumo\_\allowbreak{}360} & TEA de préstamos de consumo hasta 360 días, empresas bancarias, MN & \% (tasa efectiva anual) & Mensual & PN07846NM \\
\texttt{tea\_\allowbreak{}bcos\_\allowbreak{}consumo\_\allowbreak{}mas360} & TEA de préstamos de consumo a más de 360 días, empresas bancarias, MN & \% (tasa efectiva anual) & Mensual & PN07847NM \\
\texttt{tea\_\allowbreak{}cmac\_\allowbreak{}pyme\_\allowbreak{}360} & TEA de préstamos a pequeña y microempresa hasta 360 días, cajas municipales, MN & \% (tasa efectiva anual) & Mensual & PN07861NM \\
\texttt{tea\_\allowbreak{}cmac\_\allowbreak{}pyme\_\allowbreak{}mas360} & TEA de préstamos a pequeña y microempresa a más de 360 días, cajas municipales, MN & \% (tasa efectiva anual) & Mensual & PN07862NM \\
\texttt{tea\_\allowbreak{}crac\_\allowbreak{}pyme\_\allowbreak{}360} & TEA de préstamos a pequeña y microempresa hasta 360 días, cajas rurales, MN & \% (tasa efectiva anual) & Mensual & PN07879NM \\
\texttt{tea\_\allowbreak{}crac\_\allowbreak{}pyme\_\allowbreak{}mas360} & TEA de préstamos a pequeña y microempresa a más de 360 días, cajas rurales, MN & \% (tasa efectiva anual) & Mensual & PN07880NM \\
\texttt{ipc} & Índice de precios al consumidor de Lima Metropolitana (dic. 2021 = 100) & índice & Mensual & PN38705PM \\
\bottomrule
\end{longtable}
}

\end{singlespace}


\end{document}
